\vspace{4pt}\noindent\textbf{Embedding models.} TransE~\cite{Bordes2013TranslatingEF} models a relation as a translation in embedding space; DistMult~\cite{yang2014embedding}, ComplEx~\cite{trouillon2016complex} (complex bilinear scoring), RotatE~\cite{sun2019rotate} (rotations, able to represent symmetric and compositional patterns) and ConvE~\cite{dettmers2017convolutional} improve on it. A large re-evaluation found that classic models are competitive once their hyperparameters and training are tuned~\cite{rossi2020knowledge}, which motivates our tuning of the TransE baseline. PyKEEN~\cite{ali2020pykeen} provides such baselines and small rule-rich datasets; we reimplement TransE instead, to control the splits, the fold-in procedure and the tie policy.

\vspace{4pt}\noindent\textbf{Rules and paths.} Neural-LP~\cite{yang2017differentiable} and DRUM~\cite{sadeghian2019drum} learn soft chains of relations and are inductive. GraIL~\cite{teru2019inductive} classifies the subgraph enclosing a candidate triple and introduced inductive splits in which training and test graphs share no entities. NBFNet~\cite{zhu2021neural} generalises path-based reasoning to a Bellman-Ford style message-passing network conditioned on the query; our Path-MP is a reduced version of this idea (fixed depth, sum aggregation, no learned path pruning). Relational message passing~\cite{wang2020relational} and its fully inductive variant~\cite{geng2022relational} pass messages over relations rather than entities. R-GCN~\cite{schlichtkrull2017modeling} mixes message passing with entity-specific parameters and is thus still transductive. ULTRA~\cite{galkin2023foundation} pre-trains one model that transfers to unseen graphs and relation vocabularies. Open problems for the inductive setting, including benchmark design, are discussed by Galkin et al.~\cite{galkin2022open}.

\vspace{4pt}\noindent\textbf{What is missing.} The inductive benchmarks above are derived from existing graphs whose generative rules are unknown; the mix of rule-inferable and non-inferable test triples is unknown as well. Our contribution is the opposite trade: a toy world in which each relation is exactly rule-derived, so that differences between models, and between the two settings, can be attributed to model properties and observed-graph density rather than to unknown benchmark structure. This comes at the cost of realism, which we discuss in Section~\ref{sec:limitations}.
